(1) 两个等质量点质体各为 \(m\)，以质量可忽略的刚杆相连，距质心各为 \(l\)，绕中心天体质量 \(M\) 的圆轨道质心半径 \(R\) 运动；\(l\ll R\)，轨道角速率 \(n=\sqrt{GM/R^3}\) 由外部轨道约束维持。以杆与指向天体的局部径向的夹角 \(\theta\) 为相对姿态角，保留引力势到 \(O(l^2/R^2)\)，推导重力梯度力矩和相对姿态运动方程。
(2) 构造轨道共转系中的姿态积分，完整区分径向、切向的平衡稳定性，求径向锁定的小振动频率及分界能量。
(3) 从静止转角 \(0<\theta_m<\pi/2\) 释放，使用完整椭圆积分 \(K(k)\) 求有限振幅周期，给出小振幅修正以及趋近分界姿态时的对数发散系数。
(4) 用实验室系的自转角动量核对重力梯度力矩；指出惯性系自转能为何不单独守恒，写出共转Jacobi型守恒量及 \(l/R\to0\) 时四极展开的边界。
